\subsubsection{Espacio de trayectorias y medida invariante entre hojas}
\label{sec:medida-retorno-hojas}

La topología toroidal de APP fija el soporte de las traslaciones cardinales
y sus clases de enrollamiento. La correspondencia entre hojas procede de
otra operación sobre ese mismo soporte: la selección temporal de la
evaluación aditiva o multiplicativa. Su tipado geométrico es
\[
 \eta_{\rm av}(\Sigma)=\mathscr H_{\rm ar},\qquad
 \eta_{\rm av}(\Pi)=\mathscr H_{\rm vol}.
\]
La primera hoja corresponde a la lectura aditiva de frontera; la segunda,
a la lectura multiplicativa de interior. Este mapa transporta las etiquetas
de modo y conserva la incidencia \eqref{eq:fav}. El cuadrado de la razón
de autoescala se determina mediante las trayectorias de esa incidencia.

\Needspace{8\baselineskip}
\begin{definition}[Envolvente simbólico de memoria uno]
Sea
\[
 X_{\rm av}=\{(s_k)_{k\in\mathbb Z}\in
 \{\mathscr H_{\rm ar},\mathscr H_{\rm vol}\}^{\mathbb Z}:
 (F_{\rm av})_{s_k,s_{k+1}}=1\text{ para todo }k\}.
\]
El desplazamiento de índices actúa sobre este espacio. Es el menor sistema
simbólico de memoria uno que contiene las sucesiones modales y conserva
sus pares consecutivos: admite las dos transiciones cruzadas y la
autorretención volumétrica. La fase, la orientación cardinal y los registros
de la trayectoria permanecen en el estado TPK del que se obtiene este
envolvente.
\end{definition}

\begin{proposition}[Ponderación de las trayectorias entre hojas]
En el orden areal--volumétrico, la medida invariante de entropía máxima
de \(X_{\rm av}\) tiene matriz de transición y distribución estacionaria
\begin{equation}
 P_{\rm av}=\begin{pmatrix}0&1\\\varphi^{-2}&\varphi^{-1}\end{pmatrix},
 \qquad
 (\mu_{\rm ar},\mu_{\rm vol})
 =\frac{(1,\varphi^2)}{1+\varphi^2}.
 \label{eq:medida-hojas-exacta}
\end{equation}
En particular, \(\mu_{\rm ar}/\mu_{\rm vol}=\varphi^{-2}\).
\end{proposition}
\begin{proof}
La incidencia ya construida tiene vector propio positivo
\(r=(1,\varphi)^{\mathsf T}\) y autovalor \(\varphi\). La
normalización
\[
 (P_{\rm av})_{ij}
 =\frac{(F_{\rm av})_{ij}r_j}{\varphi r_i}
\]
da la matriz indicada. Sus filas suman uno por
\(\varphi^{-2}+\varphi^{-1}=1\). La simetría de \(F_{\rm av}\)
implica que los pesos proporcionales a \(r_i^2\) satisfacen balance
detallado; por tanto, \(\mu P_{\rm av}=\mu\). La irreducibilidad y la
autorretención volumétrica determinan una única distribución estacionaria.

Para verificar la maximalidad, en toda transición permitida se cumple
\(\log(P_{\rm av})_{ij}=\log r_j-\log r_i-\log\varphi\).
Bajo cualquier medida invariante, los términos de extremo se cancelan
en promedio. Si \(\nu_i\) son sus pesos de vértice y
\(q_i(j)=\nu(s_1=j\mid s_0=i)\), para \(\nu_i>0\), entonces
\[
 h_\nu\leq H_\nu(s_1\mid s_0)
 =\log\varphi-\sum_{i:\nu_i>0}\nu_iD(q_i\Vert P_i)
 \leq\log\varphi,
\]
donde \(D(q\Vert p)=\sum_jq_j\log(q_j/p_j)\) es la entropía
relativa en los pares permitidos y se adopta \(0\log0=0\).
La medida markoviana de \eqref{eq:medida-hojas-exacta} alcanza la cota.
La igualdad en la primera desigualdad exige la propiedad de Markov;
la igualdad en la segunda exige \(q_i=P_i\). La estacionariedad y
la irreducibilidad fijan entonces \(\mu\), lo que establece la
unicidad. Esta es la medida de Parry del envolvente definido.
\end{proof}

La frecuencia \((1/3,2/3)\) cuenta los instantes del calendario primitivo;
\(\mu\) pondera las trayectorias del envolvente. Su razón cuadrática
expresa el producto de las amplitudes de incidencia de entrada y salida.
La memoria y las restricciones de fase del TPK conservan su dominio
propio; la entropía del envolvente corresponde al sistema simbólico
especificado en esta definición.

\subsubsection{Operador de clausura y límite de los retornos marcados}
\label{sec:operador-retorno-marcado}

La coordenada de clausura \(\pi\), generada por el lector regional,
interviene después de la construcción de la incidencia modal. Sean
\(e_{\rm ar}=(1,0)^{\mathsf T}\),
\(e_{\rm vol}=(0,1)^{\mathsf T}\) y
\(P_i=e_ie_i^{\mathsf T}\) los proyectores de hoja. Las condiciones
del lector son conservación de ambas hojas, normalización volumétrica
unitaria y marca areal igual a la coordenada de clausura. Estas condiciones
determinan el operador positivo
\begin{equation}
 D_\pi=\pi P_{\rm ar}+P_{\rm vol}
       =\begin{pmatrix}\pi&0\\0&1\end{pmatrix}.
 \label{eq:operador-marca-pi}
\end{equation}
La diagonalidad expresa la conservación de hojas y las dos entradas
expresan sus normalizaciones; de este modo queda explícita la unicidad
del operador para el lector fijado.

\Needspace{8\baselineskip}
\begin{proposition}[Razón de retornos de frontera]
Para \(n\geq1\), la evaluación marcada
\begin{equation}
 \Theta_n=
 \frac{\langle e_{\rm ar},D_\pi F_{\rm av}^{\,n}e_{\rm ar}\rangle}
 {\langle e_{\rm vol},F_{\rm av}^{\,n}e_{\rm vol}\rangle}
 =\pi\frac{F_{n-1}}{F_{n+1}}
 \label{eq:retorno-marcado-finito}
\end{equation}
satisface
\begin{equation}
 \lim_{n\to\infty}\Theta_n
 =\pi\frac{\mu_{\rm ar}}{\mu_{\rm vol}}
 =\frac{\pi}{\varphi^2}.
 \label{eq:retorno-marcado-limite}
\end{equation}
\end{proposition}
\begin{proof}
Las entradas diagonales de \(F_{\rm av}^{\,n}\) cuentan los caminos
de longitud \(n\) que regresan a su hoja inicial. La fórmula de potencias
ya demostrada proporciona \(F_{n-1}\) y \(F_{n+1}\). El operador
\(D_\pi\) multiplica por \(\pi\) el primer recuento. Finalmente,
\(F_{n-1}/F_{n+1}\) es el producto de dos cocientes consecutivos
que convergen a \(\varphi^{-1}\); su límite es
\(\varphi^{-2}\). La identidad con la razón de pesos se sigue de
\eqref{eq:medida-hojas-exacta}.
\end{proof}

El operador de clausura aparece una sola vez, en la evaluación terminal
del retorno. La transferencia \(F_{\rm av}\) permanece determinada
por el calendario. Así se compone la información de dos lectores del
mismo estado: la incidencia modal determina la ponderación cuadrática
y la región de clausura aporta su coordenada a la frontera.

La composición conserva también el refinamiento. Para cilindros positivos
\(I=[a,b]\) y \(J=[c,d]\), la imagen del lector
\(\mathcal R(x,y)=x/y^2\) es el intervalo
\[
 \mathcal Q(I,J)=\left[\frac{a}{d^2},\frac{b}{c^2}\right].
\]
La monotonía en cada argumento prueba que los cilindros encajados de
clausura y autoescala producen intervalos encajados de la razón; la
positividad de sus límites y la contracción de sus diámetros determinan
\(\pi/\varphi^2\). La expresión escalar retiene así una cadena
de operaciones compatible con la profundidad, que se aplica a
continuación en sus dos orientaciones.
